\section{Lecture 1: Foundations of Deep Learning}

\subsection{Motivation and Application Areas}
Deep learning is a family of machine learning methods used in many master-level
applications:
\begin{itemize}
    \item Computer Vision
    \item Natural Language Processing (NLP)
    \item Speech and audio processing
    \item Games and intelligent agents
\end{itemize}

The key idea is to learn a function from data rather than hand-crafting rules.

\subsection{Supervised Learning Setup}
In supervised learning, we observe a dataset
\[
\mathcal{D} = \{(\mathbf{x}_1, y_1), (\mathbf{x}_2, y_2), \dots, (\mathbf{x}_I, y_I)\},
\]
where $\mathbf{x}_i$ is an input and $y_i$ is the target label/value.

We aim to learn a function
\[
f:\mathbf{x}\mapsto y.
\]

Examples:
\begin{itemize}
    \item Scalar input example: $x=$ height, $y=$ age category.
    \item Medical imaging example:
    \[
    \mathbf{x} = \text{MRI image vector} = [200,310,\dots,0], \quad
    y \in \{\text{Alzheimer: yes},\text{Alzheimer: no}\}.
    \]
\end{itemize}

For each training example, the model prediction is
\[
\hat{y}_i = f(\mathbf{x}_i;\boldsymbol{\phi}),
\]
where $\boldsymbol{\phi}$ are model parameters.

\subsection{Model Family and Linear Regression}
We do not search over all possible functions; we choose a hypothesis/model family.
A simple family is linear regression:
\[
f(x;\boldsymbol{\phi}) = \phi_0 + \phi_1 x.
\]
For classification, $f(\mathbf{x}_0;\boldsymbol{\phi})$ can be interpreted (after a
suitable output transform) as the probability of a class, e.g., Alzheimer diagnosis.

\subsection{Training Objective and Loss Function}
To select the best parameters, we minimize a loss on the dataset.
For squared error:
\[
L(\mathcal{D},\boldsymbol{\phi}) =
\sum_{i=1}^{I}\left(f(\mathbf{x}_i;\boldsymbol{\phi}) - y_i\right)^2.
\]
This is the \emph{sum of squares} criterion.

\subsection{Gradient-Based Optimization}
Parameters are updated with gradient descent:
\[
\boldsymbol{\phi} \leftarrow \boldsymbol{\phi} - \alpha \nabla_{\boldsymbol{\phi}}L(\boldsymbol{\phi}),
\]
where $\alpha > 0$ is the learning rate.

Directional derivative intuition:
for a unit direction vector $\mathbf{v}$, the rate of change is
\[
\mathbf{v}^{\top}\nabla_{\boldsymbol{\phi}}L.
\]
If the angle $\theta$ between $\mathbf{v}$ and $\nabla L$ satisfies
$\cos\theta=1$, change is maximally increasing; if $\cos\theta=-1$, change is
maximally decreasing.

\subsection{From Shallow Neural Networks to Piecewise Linear Models}
A shallow neural network with one hidden layer and ReLU activations:
\[
h_j = a(z_j), \quad z_j = \theta_{j0}+\theta_{j1}x,\quad
a(z)=\max(0,z).
\]

With two hidden units:
\[
f(x;\boldsymbol{\phi}) = \phi_0 + \phi_1 h_1 + \phi_2 h_2
= \phi_0 + \phi_1 a(\theta_{10}+\theta_{11}x) + \phi_2 a(\theta_{20}+\theta_{21}x).
\]

One possible parameter set is
\[
\boldsymbol{\phi}=
\{\phi_0,\phi_1,\phi_2,\theta_{10},\theta_{11},\theta_{20},\theta_{21}\}.
\]

Each hidden unit introduces a kink at $\theta_{j0}+\theta_{j1}x=0$.
With $D$ hidden units in 1D input, the model can form up to $D+1$ linear regions.

\newpage

\subsection{Diagram: Accumulation of ReLU-Based Linear Pieces}
\begin{figure}[h!]
\centering
\begin{tikzpicture}[scale=0.95]
    % Axes
    \draw[->] (-0.2,0) -- (8.2,0) node[right] {$x$};
    \draw[->] (0,-0.2) -- (0,4.6) node[above] {$f(x)$};

    % Combined piecewise-linear function with 3 joints (4 linear regions)
    \draw[very thick,red]
        (0,0.6) -- (1.8,0.6) -- (3.4,1.3) -- (5.2,2.6) -- (7.6,3.0);

    % Joint markers
    \filldraw[black] (1.8,0.6) circle (1.5pt);
    \filldraw[black] (3.4,1.3) circle (1.5pt);
    \filldraw[black] (5.2,2.6) circle (1.5pt);

    % Vertical guides for shifted breakpoints
    \draw[dashed] (1.8,0) -- (1.8,0.75);
    \draw[dashed] (3.4,0) -- (3.4,1.45);
    \draw[dashed] (5.2,0) -- (5.2,2.75);

    % Labels for shift locations
    \node at (1.8,-0.25) {\scriptsize $s_1$};
    \node at (3.4,-0.25) {\scriptsize $s_2$};
    \node at (5.2,-0.25) {\scriptsize $s_3$};

    % Region and model label
    \node[red] at (5.8,3.45) {\scriptsize $f(x)=\phi_0+\sum_{j=1}^{3}\phi_j a(x-s_j)$};
    \node at (6.6,0.45) {\scriptsize 4 linear regions};
\end{tikzpicture}
\caption{Combination of three shifted linear components producing three joints (kinks) in the final piecewise-linear function.}
\end{figure}

\subsection{Diagram: Input-Hidden-Output Architecture}
\begin{figure}[h!]
\centering
\begin{tikzpicture}[
    x=1cm, y=1cm,
    neuron/.style={circle, draw, minimum size=9mm},
    >=Latex
]
    % Nodes
    \node[neuron] (x1) at (0,0) {$x$};
    \node[neuron] (h1) at (3,1.6) {$h_1$};
    \node[neuron] (h2) at (3,0.0) {$h_2$};
    \node[neuron] (h3) at (3,-1.6) {$h_3$};
    \node[neuron] (y)  at (6,0.0) {$y$};

    % Edges input -> hidden
    \draw[->] (x1) -- (h1);
    \draw[->] (x1) -- (h2);
    \draw[->] (x1) -- (h3);

    % Edges hidden -> output
    \draw[->] (h1) -- (y);
    \draw[->] (h2) -- (y);
    \draw[->] (h3) -- (y);

    % Labels
    \node at (1.6,2.2) {\scriptsize $z_j=\theta_{j0}+\theta_{j1}x$};
    \node at (1.6,-2.2) {\scriptsize $h_j=a(z_j)$};
    \node at (4.7,-2.2) {\scriptsize $f(x)=\phi_0+\sum_{j=1}^{3}\phi_j h_j$};
\end{tikzpicture}
\caption{Shallow network with one input, three hidden ReLU units, and one output.}
\end{figure}
